\section{A uniform isotropic spectral gap}\label{sec:isotropic-gap}

On a smooth closed oriented homotopy four-sphere let \([g]\) be
a positive Yamabe class and take its volume-one Yamabe metric
\(g\). Put \(e=E_+(g)=E_-(g)\), \(Y=Y([g])\),
\(D=\int |Z_g|^2dV_g\), and
\[
 L=\max\{\lambda_{\max}(W^+),\lambda_{\max}(W^-)\},\qquad
 \sigma=\int L^2dV,\qquad \mathcal P([g])=4e+3\sigma .
\]
The integrals defining \(\mathcal P\) are conformally invariant.
Gauss--Bonnet and the trace-free rank-three spectral bounds give
\begin{equation}\label{eq:isotropic-quantity}
 \mathcal P-8T=D+\frac{36\sigma-Y^2}{12},\qquad
 \frac92e\le\mathcal P\le8e .
\end{equation}
In fact, for a block with largest eigenvalue \(z\), its squared
Hilbert--Schmidt norm lies between \(3z^2/2\) and \(6z^2\).
The lower bound uses Cauchy--Schwarz on the two smaller
eigenvalues whose sum is \(-z\). The upper bound uses that each
of those eigenvalues is at most \(z\), with extreme pair \((-2z,z)\).
The bounds imply \(e/6\le\sigma\le4e/3\) by signature.

\begin{theorem}[Recognition with a spectral term]\label{thm:isotropic-gap}
There is a universal \(\kappa>0\) such that every positive Yamabe
class on a nonstandard smooth homotopy four-sphere has
\(\mathcal P\ge8T+\kappa\). For any fixed \(0<\varepsilon<\kappa\),
the bound \(\inf_{Y([g])>0}\mathcal P([g])\le8T+\varepsilon\)
recognizes the standard smooth type, including an unattained
infimum. An empty family has infinite infimum.
\end{theorem}

\begin{proof}
On a nonstandard type \(36\sigma-Y^2>0\) by the strict and
equality classifications \cite[Theorems~1.4 and 1.6]{ChenZhu2014Conformal}.
If the uniform gap failed, choose positive classes on possibly
varying nonstandard types with \(\mathcal P_i\to8T\).
Equation~\eqref{eq:isotropic-quantity} gives
\(D_i\to0\), \(36\sigma_i-Y_i^2\to0\), and
\(\limsup e_i\le16T/9\). The sources are eventually in a fixed
strict band \(e_i\le2T-\delta\) for some \(\delta>0\).
In that band Gauss--Bonnet and the Yamabe upper bound give
\(\sqrt{48\delta}\le Y_i\le8\sqrt6\,\pi\). Pass to a subsequence
with \(Y_i\to y>0\), so that \(6\sqrt{\sigma_i}\to y\).

Take the volume-one modified Yamabe representatives \(q_i\)
as in Section~\ref{sec:endpoint-charge}. They satisfy
\[
 R_{q_i}-6L_{q_i}=c_i\le0,\qquad
 Y_i\le c_i+6\int L_{q_i}dV_{q_i}\le c_i+6\sqrt{\sigma_i}
 \quad\cite[Lemma~2.1 and Theorem~1.7]{ChenZhu2014Conformal}.
\]
Thus \(c_i\to0\) and \(\int L_{q_i}dV_{q_i}-\sqrt{\sigma_i}\to0\).
Since the volume is one, \(L_{q_i}\to y/6\) and \(R_{q_i}\to y\)
in \(L^2(q_i)\). In this gauge Gauss--Bonnet gives
\[
 \int |Z_{q_i}|^2dV_{q_i}
 =D_i+\frac{\int R_{q_i}^2dV_{q_i}-Y_i^2}{12}\longrightarrow0.
\]
Smooth approximations of the conformal factors in \(C^2\),
chosen separately on each source and rescaled to volume one,
give smooth representatives \(h_i\) with these same limits and
\(R_{h_i}-6L_{h_i}=o(1)\) uniformly. In particular \(R_{h_i}\ge-1\)
eventually.

Write \(U_i=|\Rm_{h_i}|_{\rm tens}\) for the full ordered-index
norm. The upper spectral bound and the curvature decomposition give
\[
 U_i\le\sqrt{48}\,L_{h_i}+\sqrt2\,|Z_{h_i}|+|R_{h_i}|/\sqrt6.
\]
For a fixed \(K_0>(\sqrt{48}/6+1/\sqrt6)8\sqrt6\,\pi\) this
implies \(\|(U_i-K_0)_+\|_{2,h_i}\to0\). The Yamabe inequality
gives a uniform Sobolev constant since
\[
 (Y_i-\|R_{h_i}-Y_i\|_2)\|f\|_4^2
 \le6\|df\|_2^2+Y_i\|f\|_2^2.
\]
The initial Einstein defects \(\|\Ric_{h_i}-(y/4)h_i\|_2\)
tend to zero as well. Apply Lemma~\ref{lem:endpoint-smoothing}.
The unnormalized flows exist on a common short interval,
and at any fixed interior positive time a subsequence converges
smoothly to a positive Einstein metric on a closed manifold,
via whole-manifold diffeomorphisms. We use also the lemma's
uniform Sobolev estimate along the flow and the estimate
\(\sup_t\|v_i(t)\|_2\to0\) for its evolving excess
\(v_i=(|\Rm_{h_i(t)}|_{\rm tens}-\phi(t))_+\) with bounded \(\phi\).

For either diagonal chiral curvature-operator block
\(A=W^\pm+RI/12\) let \(a\le b\le d\) be its eigenvalues.
Set \(F_\pm=a+b=R/6-\lambda_{\max}(W^\pm)\). Nonnegative
isotropic curvature is equivalent to \(F_+,F_-\ge0\)
\cite[Section~3, (3.1)--(3.2)]{ChenZhu2014Conformal}.
In moving orthonormal frames Hamilton's block equation is
\((D_t-\Delta)A=2(A^2+2\operatorname{adj}A+BB^*)\), using
the appropriate off-diagonal block or its transpose \(B\)
\cite[Section~2, pp.~157--158]{Hamilton1986}.
The trace on any minimizing two-plane gives the lower heat
bound \(2(a^2+b^2)+4dF_\pm\) for \(F_\pm\) after discarding
the nonnegative off-diagonal term. Hence \(w_\pm=(-F_\pm)_+\)
obeys \((\partial_t-\Delta)w_\pm\le C|\Rm|_{\rm tens}w_\pm\)
weakly. One may smooth \(w_\pm\) as the convex maximum of zero
and the negative two-plane traces in the symmetric endomorphism
fibers. Its nonnegative Hessian has the required Laplacian sign.
The active traces have the asserted bound also at zero, and
the quadratic reaction error tends to zero with the mollifier
scale on each smooth flow.

Initially \(w_\pm=o(1)\) uniformly by the modified scalar
equation. Testing its inequality by \(w_\pm\) using \(R\ge-1\)
and \(|\Rm|_{\rm tens}\le v_i+\phi\) gives
\[
 \frac d{dt}\|w_\pm\|_2^2+2\|dw_\pm\|_2^2
 \le C\|w_\pm\|_2^2+C\|v_i\|_2\|w_\pm\|_4^2.
\]
The uniform Sobolev constant absorbs the last term for large
\(i\), and Gronwall gives \(\sup_t\|w_\pm(t)\|_2\to0\).
Consequently the smooth closed positive-time limit has
nonnegative isotropic curvature.

The positive Einstein limit \(h_\infty\) realizes its class Yamabe
constant. Indeed, if \(h_\infty=u^{-2}\widehat g\) for a Yamabe
minimizer \(\widehat g\), the trace-free Ricci transformation gives
\(u Z_{\widehat g}=-2(\nabla^2_{\widehat g}u)^0\). Integrating its
pairing with the divergence-free \(Z_{\widehat g}\) makes \(\widehat g\)
Einstein. The remaining Hessian is pure trace, so either \(u\) is
constant or Obata's theorem gives the standard round type directly
\cite{Obata1962Sphere}. In the constant case the homothety
preserves the normalized scalar integral. Thus
\[
 Y([h_\infty])^2=\int R_{h_\infty}^2dV_{h_\infty}
 \ge36\int L_{h_\infty}^2dV_{h_\infty}.
\]
At strict inequality Chen--Zhu's classification gives the
standard smooth sphere in the simply connected case. At equality
the extra cases have nonzero second homology, since they are
\(\mathbb{CP}^2\) or have a compact simply connected cover
conformal to a product of two surfaces
\cite[Theorems~1.4 and 1.6]{ChenZhu2014Conformal}. The limit is
therefore standard, as are all sufficiently late sources of the
subsequence. This contradiction gives the universal gap. If
the infimum on any type is at most \(8T+\varepsilon<8T+\kappa\),
an actual class has energy quantity strictly below \(8T+\kappa\)
by the definition of infimum.
\end{proof}

In the notation \(B=4e/3-\sigma\) of Section~\ref{sec:endpoint-charge},
the theorem gives \(8(e-T)\ge3B+\kappa\) on a nonstandard type.
The theorem supplies a non-effective lower bound for
\(D+(36\sigma-Y^2)/12\) throughout the strict band and also for
arbitrary positive classes above the band. The upper bound
\(\mathcal P\le8e\) recovers a uniform recognition interval for
\(e\) alone.

\begin{corollary}[Where a vanishing spectral deficit can persist]
\label{cor:isotropic-ricci-tail}
Suppose \(e_i\le2T-\delta\) with \(\delta>0\) for positive Yamabe
classes on nonstandard types and \(36\sigma_i-Y_i^2\to0\).
Write \(\mu_i=dV_{q_i}\) for the volume measure of the volume-one
modified metric. On a subsequence there are measurable sets \(H_i\) with
\[
 \mu_i(H_i)\to0,\qquad
 \liminf_i\int_{H_i}|Z_{q_i}|^2d\mu_i>0,\qquad
 \int_{H_i}(R_{q_i}^2+|W_{q_i}|_{\rm tens}^2)d\mu_i\to0 .
\]
\end{corollary}
\begin{proof}
Suppose instead the high tails of \(|Z_{q_i}|\) are uniformly
integrable in squared norm, that is,
\(\lim_{K\to\infty}\limsup_i\|(|Z_{q_i}|-K)_+\|_2=0\).
As in the theorem, the spectral deficit gives \(c_i\to0\) as real
numbers, and both \(R_{q_i}-Y_i\) and \(L_{q_i}-Y_i/6\) tend to zero
in \(L^2(q_i)\).
The curvature decomposition therefore gives the same tail condition
for the full curvature norm of the smooth approximations.
We explain why one sufficiently small tail at a sufficiently large
fixed cutoff suffices here. In the proof of
Lemma~\ref{lem:endpoint-smoothing} keep the two coefficients of the
ordinary Sobolev inequality separate. With \(K\ge1\) and
\(\phi\le2K\), absorption in \eqref{eq:smoothing-weighted-sob}
gives a gradient coefficient bounded independently of \(K\) and
a lower-order coefficient at most \(C(1+K)\), as long as
\(\|v_i\|_2\) is below a fixed small constant. In the energy
estimates for \(v_i\) and \(w_\pm\) the gradient absorption
threshold is thus independent of \(K\). The remaining Gronwall
coefficient is at most \(C(1+K)\). On a horizon \(c/(1+K)\)
the growth factor is uniform. The ball cutoff iteration on radii
\(r\le(1+K)^{-1/2}\) gives a uniform volume ratio floor since
\((1+K)r^2\le1\). A curvature record above a sufficiently large
multiple of \(1+K\) uses only these radii, so its lower bound on
the squared excess is also independent of \(K\).
Choose \(K\) once by the tail condition and take late indices.
The common flow and smooth positive-time compactness follow,
and the negative chiral sums still tend to zero in squared norm.
The limit has nonnegative isotropic curvature.

For completeness, a smooth closed homotopy sphere with this curvature
has a positive-isotropic metric. Restart its Ricci flow; the lower
pair-sum inequality preserves both nonnegative sums. The scalar
curvature is nonnegative and somewhere positive, since if it vanished
identically both Weyl blocks would vanish and Gauss--Bonnet would
contradict positive Euler characteristic. At positive time scalar
curvature is positive. If a lower pair sum vanishes there, the
strong minimum principle gives vanishing on all earlier positive
slices. The reaction bound forces its two eigenvalues to be zero.
The block has rank one with positive trace. For a local section of
its kernel, differentiate its annihilation by the block, and test
the block evolution on that section. The sum of the off-diagonal
square and the squared norm of the component of its derivative in
the positive line is zero. Thus the kernel and its orthogonal line
are parallel. Simple connectivity gives a nonzero parallel two-form,
contrary to \(b_2=0\). The strict positive-isotropic classification
now gives the standard sphere \cite[Main Theorem]{ChenTangZhu2012}.
This contradicts the smooth identification of the limit with
late sources.

The tail condition must therefore fail. Taking increasing cutoffs
and a subsequence gives \(H_i=\{|Z_{q_i}|>K_i\}\) with definite
squared integral and \(K_i\to\infty\). The uniformly bounded total
Ricci energy, by Gauss--Bonnet in \(q_i\), gives \(\mu_i(H_i)\to0\).
The scalar and spectral squared-integral limits and
\(|W|_{\rm tens}^2\le48L^2\) prove the last assertion.
\end{proof}
